%-------------------------------------------------------------
% CV — Luis Arce C.
% Compilar con: pdflatex cv-es.tex
%-------------------------------------------------------------
\documentclass[10pt,letterpaper]{article}

\usepackage[utf8]{inputenc}
\usepackage[T1]{fontenc}
\usepackage[spanish,es-noshorthands,es-nolists,es-noquoting]{babel}
\usepackage{tgheros}
\renewcommand{\ttdefault}{lmtt}
\renewcommand{\familydefault}{\sfdefault}
\usepackage[letterpaper,top=1.3cm,bottom=1.2cm,left=1.75cm,right=1.75cm]{geometry}
\usepackage{microtype}
\usepackage{xcolor}
\usepackage{titlesec}
\usepackage{enumitem}
\usepackage{needspace}
\usepackage[hidelinks]{hyperref}

%----------------------------- Paleta
\definecolor{ink}{HTML}{1F2328}
\definecolor{accent}{HTML}{1F3A5F}
\definecolor{muted}{HTML}{667085}
\definecolor{rule}{HTML}{D0D5DD}
\definecolor{todo}{HTML}{B42318}

\hypersetup{pdftitle={Luis Arce C. — CV},pdfauthor={Luis Arce C.}}
\pagestyle{empty}
\setlength{\parindent}{0pt}
\setlength{\parskip}{0pt}
\color{ink}
\linespread{1.03}

%----------------------------- Secciones
\titleformat{\section}
  {\color{accent}\footnotesize\bfseries}
  {}{0pt}
  {}
  [\vspace{2pt}{\color{rule}\titlerule[0.5pt]}]
\titlespacing*{\section}{0pt}{9pt}{5pt}
\newcommand{\cvsection}[1]{\needspace{5\baselineskip}\section{\MakeUppercase{#1}}}

%----------------------------- Listas
\setlist[itemize]{
  leftmargin=1.05em, labelsep=0.5em, topsep=2.5pt, itemsep=1.2pt, parsep=0pt, partopsep=0pt,
  label={\color{muted}\raisebox{0.25ex}{\tiny$\bullet$}},
  before=\small, beginpenalty=10000
}

%----------------------------- Bloques
% \cventry{Empresa}{Ubicación}{Cargo}{Periodo}
\newcommand{\cventry}[4]{%
  \vspace{6pt}%
  {\bfseries #1}\hfill{\small\color{muted}#4}\par
  \vspace{1pt}%
  {\small #3}\hfill{\small\color{muted}#2}\par
  \vspace{1pt}%
}
% \cvproject{Proyecto}{Descriptor}{Periodo}
\newcommand{\cvproject}[3]{%
  \vspace{5pt}%
  {\small{\color{accent}\bfseries #1}\quad{\color{muted}#2}\hfill{\color{muted}#3}}\par\nobreak
}
% \cvskill{Categoría}{Contenido}
\newcommand{\cvskill}[2]{%
  {\small\textbf{#1}\quad #2}\par\vspace{2.5pt}%
}
\newcommand{\sep}{\hspace{0.55em}{\color{rule}|}\hspace{0.55em}}
\newcommand{\pendiente}[1]{{\color{todo}[#1]}}

\begin{document}

%============================= ENCABEZADO
{\fontsize{22}{26}\selectfont\bfseries Luis Arce C.}\par
\vspace{5pt}
{\color{accent}Ingeniero de Software Senior \textperiodcentered{} AI Engineer \textperiodcentered{} Arquitectura Backend}\par
\vspace{5pt}
{\small\color{muted}%
  Milagro, Ecuador\sep
  \href{mailto:luilliarcec@gmail.com}{luilliarcec@gmail.com}\sep
  +593 979 221 239\sep
  \href{https://www.linkedin.com/in/luilliarcec}{linkedin.com/in/luilliarcec}\sep
  \href{https://github.com/luilliarcec}{github.com/luilliarcec}}\par

%============================= PERFIL
\cvsection{Perfil}
{\small
Ingeniero de Sistemas con más de 8 años de experiencia y foco en ingeniería de IA aplicada: integro modelos de lenguaje en
productos en producción (asistentes conversacionales multi-LLM, transcripción y traducción de voz, extracción de documentos) y
diseño flujos agénticos en los que agentes de IA implementan bajo especificaciones, pruebas automatizadas y mi revisión.
Sobre esa base construyo y opero plataformas SaaS multi-tenant, microservicios e integraciones críticas en fintech, salud, logística
portuaria y recursos humanos para clientes de Ecuador, Estados Unidos y Alemania, con Laravel/PHP como núcleo y trabajo
complementario en Node.js, Java, .NET y Rust. Como líder técnico asumo arquitectura, seguridad, calidad y operación.
\par}

%============================= EXPERIENCIA
\cvsection{Experiencia profesional}

\cventry{EFY}{Remoto}{AI Engineer}{Febrero 2026 -- Actualidad}
{\small Fintech B2B y B2C para Latinoamérica: wallets cripto custodiadas, cuentas bancarias en USD, tarjetas prepago, compra y
venta de cripto, transferencias entre usuarios y un back office para Cumplimiento y Tesorería.\par}

\cvproject{Ingeniería agéntica}{Desarrollo de la plataforma con agentes de IA bajo mi dirección}{}
\begin{itemize}
  \item Más de 1.400 PRs mergeados en 19 repositorios en nueve meses (backend, frontend, contratos, infraestructura y
        conocimiento de producto), implementados en su mayoría por agentes de Claude Code que orquesto: asigno cada tarea a un
        modelo según su complejidad, corro agentes en paralelo sobre worktrees aislados (Orca) y les doy acceso a Jira, Slack,
        AWS, bases de datos y navegador mediante servidores MCP.
  \item Cada cambio pasa por revisión automática y pruebas antes del merge, y por una verificación e2e en vivo tras desplegarse
        en desarrollo. Ajusto el proceso cada semana con la telemetría de los agentes: PRs aprobados en la primera revisión
        del 85 al 92--98\,\% y e2e en verde al primer intento del 73 al 87--90\,\%.
  \item Más de 150 PRs al plugin de skills del equipo (playbooks de resolución de tickets, recorridos e2e en vivo y guards de
        «diseño primero» y de plantillas PDF) y aportes a la base de conocimiento de producto que los agentes consultan antes
        de crear o resolver una tarea.
\end{itemize}

\cvproject{Backend y pagos}{Laravel 12, Node.js/Fastify y TypeScript sobre AWS}{}
\begin{itemize}
  \item Servicios de transacciones (Laravel 12, PHP 8.4), usuarios e identidad, BFF, comprobantes, extractos y notificaciones
        sobre AWS (ECS Fargate, Lambda, Aurora PostgreSQL, Cognito, S3, SES), con sagas, outbox, idempotencia y un ledger de
        partida doble.
  \item Integraciones con banca en USD (apertura de cuentas con KYC y KYB, pagos ACH y SWIFT con compensación contable en dos
        pasos), un procesador de tarjetas prepago (emisión, recargas desde cripto, disputas y extractos), custodia cripto MPC
        y una wallet en pesos colombianos. Suite e2e en vivo contra los cuatro sandboxes que valida la API, el proveedor y el
        ledger; la auditoría del conector de tarjetas dejó 12 hallazgos: 10 corregidos y 2 refutados con evidencia.
  \item Seguridad y cumplimiento: cerré accesos indebidos entre empresas (IDOR) con permisos HMAC-SHA256 de 60 segundos
        firmados por el BFF, llevé FATCA a un envío estructurado e inmutable e implementé en cinco repositorios la máquina de
        estados del alta de empresas. Bajé la suite de 1.648 pruebas de 192 s a 38 s (6 s en paralelo).
\end{itemize}

\cvproject{Frontend y diseño}{Next.js 16, Expo/React Native y Claude Design}{}
\begin{itemize}
  \item Casi 300 PRs en el monorepo de front (pnpm y Turborepo): apps web B2C y B2B en Next.js 16 y app móvil en Expo,
        conectadas al BFF por un contrato OpenAPI con mocks (MSW) y control de compatibilidad (oasdiff).
  \item Flujo de diseño a código: las pantallas se diseñan en Claude Design y llegan al repositorio como drops versionados que
        se portan a Next.js y React Native con un sello de integridad por pantalla. Implementé el gate de CI «diseño primero»,
        que bloquea toda pantalla B2C que no salga de ese flujo, y gestiono los cambios de diseño como change requests versionados.
  \item Comprobantes y extractos en PDF (react-pdf) fieles a las plantillas de Claude Design, y módulos del back office para
        Cumplimiento y Tesorería: colas de fondeo, rechazo masivo de documentos KYC/KYB, disputas e invalidación de FATCA.
\end{itemize}

\cventry{TeamQ Developers}{Quito, Ecuador \textperiodcentered{} Remoto}{Senior Fullstack Developer \textperiodcentered{} Líder Técnico}{Enero 2023 -- Actualidad}

\cvproject{RUDY}{SaaS B2B de certificados laborales para el mercado alemán \textperiodcentered{} Líder técnico}{2023 -- Actualidad}
\begin{itemize}
  \item Lideré la reescritura completa de la plataforma: tres aplicaciones heredadas (CakePHP, Laravel y Angular sobre MySQL)
        unificadas en una API Laravel 13 con Octane/FrankenPHP sobre PostgreSQL 17 y un front Next.js 16. Salió a producción
        unas 12 semanas después del primer commit, con la migración de 48 empresas y una ventana de corte de 22 minutos.
  \item Diseñé el aislamiento entre empresas con Row-Level Security de PostgreSQL, cifrado AES-256 por campo con una clave por
        empresa y rotación sin interrupción, SSO OIDC por cliente (secreto o mTLS) y doble factor de autenticación.
  \item Entrega continua con despliegue blue-green: más de 110 versiones en producción en dos meses, despliegues de menos de
        un minuto y rollback en 13 segundos. El API se sostiene con 4.801 pruebas automatizadas (97,3\,\% de cobertura),
        análisis estático en nivel 8 y especificaciones versionadas junto al código (más de 500 requisitos).
  \item Operación y seguridad: respaldos cifrados con copia externa inmutable, reconstrucción completa ensayada (RPO y RTO
        de 4 horas), runbooks de despliegue y SDLC alineado con OWASP SAMM. Resueltos los 10 hallazgos de severidad alta de
        una evaluación de seguridad externa.
\end{itemize}

\cvproject{HealthHero}{Vacunación escolar en EE.\,UU. \textperiodcentered{} Arquitecto y desarrollador principal del backend}{2023 -- Actualidad}
\begin{itemize}
  \item \textbf{ERIS}: plataforma multi-tenant para socios regionales (Laravel, MySQL en AWS RDS, SQS, S3). Autor principal del
        backend: pacientes, consentimientos, inventario de vacunas con trazabilidad por eventos, y facturación y reclamaciones
        a aseguradoras con ClaimMD (lotes, conciliación de remesas y rechazos).
  \item Integración HL7 (mensajes VXU) con los registros de inmunización (IIS) de 10 estados: una Abstract Factory resuelve la
        implementación de cada estado y su transporte (REST, SOAP o SFTP), con acuses asíncronos mediante colas diferidas.
  \item Digitalización masiva de consentimientos en papel con Azure Document Intelligence (modelos a medida y detección de la
        versión del formulario por código QR) y búsqueda histórica de pacientes sobre MongoDB y Typesense, con puntuación
        ponderada por atributo y distancia de Levenshtein.
  \item \textbf{Kronos} (Laravel 12, PHP 8.4): planificación de clínicas por temporada, coordinación de equipos de enfermería,
        nómina y optimización de rutas con Google Route Optimization. Diseñé su integración con ERIS, que actúa como autoridad
        OAuth 2.0 del ecosistema. Ambos backends suman más de 2.800 pruebas automatizadas.
\end{itemize}

\cvproject{TeamQ SBS}{Suite SaaS de microservicios para pymes \textperiodcentered{} Arquitecto y líder técnico de Omnichat}{2023 -- Actualidad}
\begin{itemize}
  \item Creé el servicio central de organizaciones (multi-tenant) y lo rediseñé sobre Laravel 12 y Keycloak. Autor de los
        paquetes compartidos de tenencia y de autenticación (usuarios y máquina a máquina) y de la plantilla base de la que
        derivan cuatro microservicios; sincronización de identidades Keycloak $\rightarrow$ RabbitMQ $\rightarrow$ servicios.
  \item \textbf{Omnichat}: bandeja omnicanal sobre WhatsApp Cloud, Messenger, Instagram y Telegram (Laravel 13, Octane, MongoDB,
        RabbitMQ), de la que soy desarrollador principal. Diseñé el motor de flujos conversacionales como máquina de estados
        (expiración por inactividad, invalidación al republicar, tope de nodos por turno contra ciclos), un renderizador
        independiente del canal y el traspaso a agentes humanos.
  \item IA multi-proveedor por organización (OpenAI, Anthropic, Gemini): cada cliente usa sus propias credenciales, aisladas
        entre peticiones en los workers de larga vida de Octane.
\end{itemize}

\cvproject{Terminal Portuario de Guayaquil (TPG)}{Sistemas corporativos e integración con Navis N4}{2023 -- 2026}
\begin{itemize}
  \item Backends Laravel con SSO Keycloak para crédito, certificados de carga y contenedores, control de acceso, calificación
        de proveedores y centro de ayuda; principal contribuidor en cuatro de ellos y autor del starter kit común.
  \item Integración con Navis N4, el sistema operativo de la terminal: primero un gateway en NestJS/Fastify (plantillas SOAP en
        MongoDB, caché en Redis) y después un middleware en Apache Camel 4 (YAML DSL, Java 21, Groovy) desplegado en
        Kubernetes con GitOps. Aporté el endpoint SOAP unificado, orquestaciones en paralelo y un worker asíncrono sobre
        Redis Sentinel.
\end{itemize}

\cvproject{TaxiApp}{API de transporte con despacho en tiempo real}{2024 -- 2026}
\begin{itemize}
  \item Construí la base del backend: seguimiento GPS en tiempo real (MongoDB y WebSockets con Laravel Reverb), push con FCM,
        chat, seis idiomas y zona horaria por usuario; después, el despacho automático con colas FIFO de conductores por geocerca.
\end{itemize}

\cvproject{Otros proyectos y aportes internos}{}{}
\begin{itemize}
  \item Carga masiva de datos en .NET 10 para una aseguradora de salud, gestión de programas para una universidad y campañas
        push para una cadena hotelera. Autor de \texttt{teamq-laravel-datatables}, paquete de código abierto con más de 10.000
        descargas en Packagist, y del starter kit Laravel de la empresa.
  \item Autor del estándar de trabajo con IA del equipo: skills de Claude Code que cubren la creación de tickets, el flujo de
        ramas y PR, y la resolución de una tarea de punta a punta (ticket, pruebas, PR, CI y despliegue).
\end{itemize}

\cventry{Codesofsas}{Guayaquil, Ecuador}{Senior Fullstack Developer}{Agosto 2020 -- Febrero 2023}
\begin{itemize}
  \item \textbf{Texcumar} (arquitecto de software): lideré el diseño y desarrollo de la aplicación de control de producción de
        nauplios de camarón (Laravel, Livewire, AlpineJS), con inventarios y reportes por etapa cubiertos por pruebas automatizadas.
  \item \textbf{Barcelona Sporting Club}: e-commerce y API para la aplicación móvil oficial del club.
  \item \textbf{ERP Durán}: desarrollo de un ERP multi-tenant con Laravel y Vue.js.
\end{itemize}

\cventry{Freelance}{Milagro, Ecuador}{Desarrollador Fullstack}{Octubre 2018 -- Actualidad}
\begin{itemize}
  \item Aplicaciones web (Laravel, Node.js) y de escritorio (C\#), y bases de datos SQL Server con lógica de negocio crítica en T-SQL.
\end{itemize}

%============================= PROYECTOS PERSONALES
\cvsection{Proyectos independientes}

\cvproject{Ferretería Calderón}{ERP para una pyme ferretera en Ecuador \textperiodcentered{} Arquitecto e ingeniero a cargo}{2026}
\begin{itemize}
  \item Punto de venta, inventario con kardex, contabilidad de partida doble y facturación electrónica ante el SRI (firma XAdES
        y envío SOAP de facturas, notas, retenciones y guías). Laravel 13, React 19 con Inertia, Octane/FrankenPHP y
        PostgreSQL 18; acceso con 2FA y passkeys.
  \item Construido con agentes de IA bajo mi dirección: especificaciones versionadas, pruebas y análisis estático como puerta de
        cada cambio (unos 100 merge requests en menos de cuatro meses) y migración del sistema anterior conciliada al centavo.
\end{itemize}

\cvproject{AI TVoice}{Traductor de voz en tiempo real para videollamadas \textperiodcentered{} Creador}{2025}
\begin{itemize}
  \item Aplicación de escritorio en Rust (Tokio, Tauri 2) que captura a la vez el audio del sistema y el micrófono, transcribe con
        Whisper y traduce con un LLM. Diseñé la cadena de audio completa: detección de voz por energía, segmentación,
        remuestreo a 16 kHz, resultados parciales en vivo y salida de voz hacia un micrófono virtual.
\end{itemize}

%============================= EDUCACIÓN
\cvsection{Educación}

\cventry{Universidad Estatal de Milagro}{Milagro, Ecuador}{Maestría en Ingeniería de Software \textperiodcentered{} Sistemas inteligentes e intensivos en datos}{Octubre 2025 -- Actualidad}
\cventry{Universidad Estatal de Milagro}{Milagro, Ecuador}{Ingeniería en Sistemas}{Febrero 2020}

%============================= HABILIDADES
\cvsection{Habilidades técnicas}
\vspace{2pt}
\cvskill{Ingeniería de IA}{Integración de LLM (OpenAI, Anthropic, Gemini), agentes y skills (Claude Code, Orca, MCP), Whisper, OCR.}
\cvskill{Backend}{PHP 8.4, Laravel (experto), Octane/FrankenPHP, Node.js (NestJS, Fastify), C\#/.NET, Java (Apache Camel), Rust.}
\cvskill{Arquitectura}{SaaS multi-tenant, microservicios, sistemas orientados a eventos, DDD, integraciones REST/SOAP/SFTP/HL7.}
\cvskill{Datos y mensajería}{PostgreSQL (RLS), MySQL, SQL Server (T-SQL), MongoDB, Redis, Typesense, RabbitMQ, AWS SQS.}
\cvskill{Seguridad e identidad}{OAuth 2.0, OIDC, Keycloak, SSO, mTLS, 2FA, cifrado por campo, gestión de secretos.}
\cvskill{Infraestructura y calidad}{Docker, AWS, Azure, GitHub Actions, despliegue blue-green, TDD (Pest, PHPUnit), PHPStan/Larastan.}
\cvskill{Frontend}{React (Inertia, Next.js), React Native (Expo), Livewire, Filament, Alpine.js, Angular, Vue.js.}
\cvskill{Idiomas}{Español (nativo), inglés (A2, en formación activa).}

\end{document}
